\begin{table}[htbp] \centering \caption{Summary Statistics Indonesia, Balanced Panel}\label{tab:summary_stats_IDN_bal} \begin{threeparttable} \begin{tabular}{l cccc} \toprule
& All & Rural & Urban & Difference \\
 &  &  &  & $t$-test \\
\midrule
Location &  &   55.2\% &   44.8\% &  \\  \addlinespace \addlinespace
Log Consumption & 12.00 & 11.82 & 12.21 & -0.39*** \\
 & (0.80) & (0.80) & (0.76) &  \\  \addlinespace
Log Income & 14.92 & 14.73 & 15.16 & -0.43*** \\
 & (1.13) & (1.14) & (1.07) &  \\  \addlinespace
Female & 0.50 & 0.48 & 0.52 & -0.04*** \\
 & (0.50) & (0.50) & (0.50) &  \\  \addlinespace
Age (years) & 43.45 & 42.59 & 44.52 & -1.93*** \\
 & (12.16) & (12.19) & (12.03) &  \\  \addlinespace
Education (years) & 7.57 & 6.63 & 8.72 & -2.09*** \\
 & (4.40) & (4.27) & (4.28) &  \\  \addlinespace
Household Size & 4.89 & 4.82 & 4.99 & -0.17*** \\
 & (2.09) & (1.98) & (2.21) &  \\  \addlinespace


\cmidrule{2-5}
Observations & 16,420 & 9,059 & 7,361 &  \\  
Individuals & 3,284 &  &  &  \\
Non-switchers &   59.6\% &  &  &  \\  \addlinespace
\bottomrule \end{tabular} \begin{tablenotes}[flushleft] \footnotesize \item{Summary statistics for Indonesia for the balanced panel across all five waves. Source: IFLS. The table reports means and standard deviations (in parentheses) based on individual-year pairs. See section 3 for further details. All variables have the same number of observations, except for income, which is missing for some observations. Income has 12,510 observations.} \end{tablenotes} \end{threeparttable} \end{table}